% =====================================================================
%  Diagrama canónico del framework como sistema de Tubos y Filtros
%  (Figura del Capítulo 3, Sección 3.2.1 de main.tex)
%
%  Source of truth del diagrama. Para regenerar el PNG que usa main.tex:
%      lualatex diagrama_arquitectura.tex
%      pdftoppm -png -singlefile -r 150 diagrama_arquitectura.pdf Diagrama-Arquitectura
%
%  El primer comando produce diagrama_arquitectura.pdf (ignorado por git);
%  el segundo produce Diagrama-Arquitectura.png (versionado, referenciado
%  por \includegraphics en main.tex).
% =====================================================================
\documentclass[tikz,border=14pt]{standalone}
\usepackage[T1]{fontenc}
\usepackage[utf8]{inputenc}
\usepackage[spanish,es-tabla,es-noindentfirst]{babel}
\usetikzlibrary{shapes.geometric,positioning,arrows.meta,calc,fit,backgrounds}

\begin{document}
\begin{tikzpicture}[
  font=\small,
  filtro/.style={draw, thick, rectangle, rounded corners=2pt,
    minimum width=2.4cm, minimum height=1.6cm, align=center},
  exterior/.style={draw, thick, cylinder, shape border rotate=90, aspect=0.25,
    minimum width=2.5cm, minimum height=1.5cm, align=center, font=\small\itshape},
  banda/.style={draw, thick, rectangle, rounded corners=3pt, align=center, fill=black!4},
  repos/.style={draw, line width=2pt, rectangle, rounded corners=3pt, align=center, fill=black!4},
  puerto/.style={fill=black, draw=black, rectangle, minimum size=5pt, inner sep=0pt},
  pnum/.style={font=\scriptsize\bfseries},
  tubo/.style={thick},
  conexion/.style={thick, dash pattern=on 4pt off 3pt},
  sint/.style={{Circle[length=2.2mm]}-{Latex[length=2.8mm,width=2.2mm]}, thick},
  memoria/.style={{Latex[length=3.4mm,width=3mm]}-{Latex[length=3.4mm,width=3mm]}, line width=1.9pt},
  sintbox/.style={draw, thick, rectangle, rounded corners=2pt, align=center,
    font=\scriptsize, fill=black!3, minimum height=0.85cm, text width=2.9cm, inner sep=3pt},
  compuesto/.style={draw, thick, dash pattern=on 3pt off 2pt, rounded corners=6pt},
]

% ---------- Filtros (banda media, y=4.5) ----------
\node[filtro] (extr)  at (3,4.5)    {Extracción};
\node[filtro] (enriq) at (6.2,4.5)  {Enriquecimiento\\semántico};
\node[filtro] (comp)  at (9.4,4.5)  {Composición\\de contexto};
\node[filtro] (invoc) at (12.6,4.5) {Invocación\\generativa};
\node[filtro] (deco)  at (15.8,4.5) {Decodificación};
\node[filtro] (norm)  at (19.0,4.5) {Normalización};

% ---------- Fuente y sumidero (separados para que la conexión punteada respire) ----------
\node[exterior] (portales) at (-1.1,4.5) {Portales\\digitales};
\node[exterior] (presenta) at (23.1,4.5) {Capa de\\presentación\\\scriptsize(tableros, reportes)};

% ---------- Puertos (cuadraditos en bordes O/E, y=4.5) ----------
\foreach \f in {extr,enriq,comp,invoc,deco,norm}{
  \node[puerto] (\f-in)  at (\f.west) {};
  \node[puerto] (\f-out) at (\f.east) {};
}

% ---------- Numeración de puertos (arriba del tubo) ----------
\node[pnum] at ($(extr-in)  +(-0.18,0.32)$) {1};
\node[pnum] at ($(extr-out) +(0.18,0.32)$)  {2};
\node[pnum] at ($(enriq-in) +(-0.18,0.32)$) {3};
\node[pnum, text=black!55] at ($(enriq-out)+(0.18,0.32)$)  {4};
\node[pnum, text=black!55] at ($(comp-in)  +(-0.18,0.32)$) {5};
\node[pnum] at ($(comp-out) +(0.18,0.32)$)  {6};
\node[pnum] at ($(invoc-in) +(-0.18,0.32)$) {7};
\node[pnum] at ($(invoc-out)+(0.18,0.32)$)  {8};
\node[pnum] at ($(deco-in)  +(-0.18,0.32)$) {9};
\node[pnum] at ($(deco-out) +(0.18,0.32)$)  {10};
\node[pnum] at ($(norm-in)  +(-0.18,0.32)$) {11};
\node[pnum] at ($(norm-out) +(0.18,0.32)$)  {12};

% ---------- Tubos tipados entre filtros (líneas sin flecha, unen puertos) ----------
\draw[tubo] (extr-out)  -- (enriq-in);
\draw[tubo, draw=black!45] (enriq-out) -- (comp-in);
\draw[tubo] (comp-out)  -- (invoc-in);
\draw[tubo] (invoc-out) -- (deco-in);
\draw[tubo] (deco-out)  -- (norm-in);

% ---------- Conexiones al entorno (punteadas, sin flecha) ----------
\draw[conexion] (portales.east) -- (extr-in);
\draw[conexion] (norm-out) -- (presenta.west);

% ---------- Sintonizadores modulares (uno por filtro que el experto configura) ----------
% Cada sintonizador es la interfaz de configuración de su propio filtro; agruparlos
% en una GUI única es decisión del implementador, no de la arquitectura.
% Enriquecimiento no expone sintonizador: su config (modelo de embedding, chunking)
% es técnica del implementador. Los parámetros de dominio del RAG (top-k, umbral)
% viven en el sintonizador de Composición.
\node[sintbox] (sint-extr) at (3,7.3)   {Selección de portales,\\reglas y prioridad};
\node[sintbox] (sint-comp) at (9.4,7.3) {Template, variables\\y parámetros de recuperación};
\node[sintbox] (sint-norm) at (19,7.3)  {Criterios de\\consolidación};
\draw[sint] (sint-extr.south) -- (extr.north);
\draw[sint] (sint-comp.south) -- (comp.north);
\draw[sint] (sint-norm.south) -- (norm.north);

% ---------- Filtro compuesto: Recuperación aumentada (Enriquecimiento + Composición) ----------
\begin{scope}[on background layer]
  \node[compuesto, fit=(enriq)(comp), inner xsep=7pt, inner ysep=12pt] (rag) {};
\end{scope}
\node[font=\small\itshape, anchor=south west] at (rag.north west) {Recuperación aumentada};
% Puertos internos del compuesto (4-5): tono atenuado para señalar su encapsulación
\node[puerto, fill=black!40, draw=black!40] at (enriq.east) {};
\node[puerto, fill=black!40, draw=black!40] at (comp.west) {};

% ---------- Banda inferior: repositorio común (borde grueso) ----------
\node[repos, minimum width=18cm, minimum height=1.2cm] (repo) at (11,0.3)
  {Repositorio común --- PostgreSQL / ORM de Odoo};
\foreach \f in {extr,enriq,comp,invoc,deco,norm}{
  \draw[memoria] (\f.south) -- (\f.south |- repo.north);
}

% ================= LEYENDA =================
\node[anchor=north west, font=\footnotesize] (leyenda) at (-1.5,-1.6) {%
  \begin{tabular}{@{}cl@{\hspace{10pt}}cl@{}}
    \tikz[baseline=-0.6ex]{\draw[thick,rounded corners=1pt] (0,-0.13) rectangle (0.5,0.13);}
      & filtro
    & \tikz[baseline=-0.6ex]{\draw[thick,fill=black] (0,-0.06) rectangle (0.12,0.06);}
      & puerto (tipado) \\[5pt]
    \tikz[baseline=-0.9ex]{%
        \draw[thick] (-0.23,0) -- (-0.23,-0.34);
        \draw[thick] (0.23,0) -- (0.23,-0.34);
        \draw[thick] (-0.23,-0.34) arc (180:360:0.23 and 0.08);
        \draw[thick] (0,0) ellipse (0.23 and 0.08);}
      & fuente / sumidero
    & \tikz[baseline=-0.6ex]{\draw[line width=1.9pt,{Latex[length=2.2mm]}-{Latex[length=2.2mm]}] (0,0) -- (0.65,0);}
      & acceso al repositorio (contenido) \\[6pt]
    \tikz[baseline=-0.6ex]{\draw[thick] (0,0) -- (0.6,0);}
      & tubo tipado (referencias)
    & \tikz[baseline=-0.6ex]{\draw[thick,{Circle[length=1.6mm]}-{Latex[length=2mm]}] (0,0) -- (0.6,0);}
      & sintonizador \\[5pt]
    \tikz[baseline=-0.6ex]{\draw[thick,dash pattern=on 4pt off 3pt] (0,0) -- (0.6,0);}
      & conexión al entorno
    & \tikz[baseline=-0.6ex]{\draw[thick,dash pattern=on 3pt off 2pt,rounded corners=2pt] (0,-0.13) rectangle (0.5,0.13);}
      & filtro compuesto \\
  \end{tabular}
};

% ================= TABLA DE PUERTOS =================
\node[anchor=north west, font=\footnotesize] (tabla) at (10.6,-1.6) {%
  \begin{tabular}{@{}rll@{\hspace{14pt}}rll@{}}
    \multicolumn{6}{@{}l}{\textbf{Referencia de puertos}\quad\footnotesize(\texttt{X[c]} denota un recordset de tipo \texttt{X} y cardinalidad $c$; el contenido reside en el repositorio común)} \\
    \multicolumn{6}{@{}l}{\footnotesize($N$: cantidad de noticias del corpus; $M$: cantidad de unidades de análisis, con $1 \le M \le N$)} \\[3pt]
    1  & Extracción (ent.)      & flujo del portal       & 7  & Invocación (ent.)     & \texttt{Prompt[M]} \\
    2  & Extracción (sal.)      & \texttt{Noticia[N]}    & 8  & Invocación (sal.)     & \texttt{RespuestaIAG[M]} \\
    3  & Enriquecimiento (ent.) & \texttt{Noticia[N]}    & 9  & Decodificación (ent.) & \texttt{RespuestaIAG[M]} \\
    4  & Enriquecimiento (sal.) & \texttt{Embedding[N]}  & 10 & Decodificación (sal.) & \texttt{Análisis[M]} \\
    5  & Composición (ent.)     & \texttt{Embedding[N]}  & 11 & Normalización (ent.)  & \texttt{Análisis[M]} \\
    6  & Composición (sal.)     & \texttt{Prompt[M]}     & 12 & Normalización (sal.)  & \texttt{Resultado} \\
  \end{tabular}
};

\end{tikzpicture}
\end{document}
